\documentclass[10pt]{article}

\usepackage[utf8]{inputenc}

\usepackage[T1]{fontenc}

\usepackage{amsmath}

\usepackage{amsfonts}

\usepackage{amssymb}

\usepackage[version=4]{mhchem}

\usepackage{stmaryrd}

\usepackage{graphicx}

\usepackage[export]{adjustbox}

\graphicspath{ {./images/} }

\usepackage{bbold}



\begin{document}

образом, изучение их строения в известной мере сводится к рассмотрению типов, к к-рым принадлежат компоненты связки, и к рассмотрению полугрупп идемпотентов (см., напр., А рхимедова полугруппа, Вполне простая полугруппа, Клиффордова полугруппа, Периодическая полугруппа, Сепаративнал полугруппа).



С. ІІ. $S_{\alpha \text { наз. к о м м у т а т и в н о й, если для соот- }}^{\text {к }}$ ветствующей конгруэнции $\rho$ факторполугруппа $S / \rho$ коммутативна; тогда $S / \rho$ - полурешетка (в этом случае часто говорят, что $S$ есть полурешетка полугрупп $S_{\alpha}$, в частности, если $S / \rho-$ цепь, то говорят, что $S$ есть цепь полугрупп $\left.S_{\alpha}\right)$. С. П. наз. п р я м о у г о л ь н о й (иногда - ма т р и чно й), если $S / \rho-$ прямоугольная полугруппа (см. Идемпотентов полугруппа). Это эквивалентно тому, что компоненты связки могут быть индексированы парами индексов $S_{i \lambda}$, где $i$ и $\lambda$ пробегают нек-рые множества $I$ и $\Lambda$ соответственно, причем для любых $S_{i \lambda}, S_{j \mu}$ выполняется включение $S_{i \lambda} S_{j \mu} \sqsubseteq S_{i \mu}$. Јюбая С. п. есть полурешетка прямоугольных связок, т. е. ее компоненты могут быть распределены на подсемейства так, что объединение компонент каждого подсемейства есть прямоугольная связка этих компонент, а исходная полугруппа разложима в полурешетку указанных объединений (т е о p eм а К л и ф ф о р д а [1]). Поскольку свойства быть полугруппой идемпотентов, полурешеткой, прямоугольной полугруппой характеризуются тождествами, на любой полугруппе $S$ для каждого из перечисленных свойств $\theta$ существует наименьшая конгруэнция, для к-рой соответствующая факторполугруппа обладает свойством $\theta$, т. е. существуют н а и б о ль ш и е (или н а и боле е д робны е) р а зложен и я $S$ в С. п., в коммутативную С. П., в прямоугольную С. П.



K специальным типам С. П. относится с и л ь н а я с в я 3 к а [4]: для любых элементов $a$ и $b$ из разных компонент произведение $a b$ равно степени одного из этих элементов. Важным частным случаем сильной связки и одновременно частным случаем цепи полугрупп является о р д ин а л ь н а я с у м м а (или пос ле дов а тельн о ант у ли р ующая с в я з а): множество ее компонент $\left\{S_{\alpha}\right\}$ линейно упорядочено и для любых $S_{\alpha}, S_{\beta}$ таких, что $S_{\alpha}<S_{\beta}$, и любых $a \in S_{\alpha}$, $b \in S_{\beta}$ имеет место равенство $a b=b a=a$. Заданием компонент и способа их упорядочения ординальная сумма определяется однозначно с точностью до изоморфизма.



Jum.: [1] C I i f f o r d A., "Proc. Amer. Math Soc $\nu, 1954$,



\begin{center}

\includegraphics[max width=\textwidth]{2023_07_08_a610a5a2eea8429141b4g-1(1)}

\end{center}



\includegraphics[max width=\textwidth, center]{2023_07_08_a610a5a2eea8429141b4g-1}

«Изв. вузов. Матем.", 1965 , №6, с. $156-65$. Л. Н. Шеврин.



СВЯЗНАЯ КОМПОНЕНТА ЕДИНИЦЫ г р у п п Ы $G$ - наибольшее связное подмножество $G^{\circ}$ топологической (или алгебраической) группы $G$, содержащее единицу этой группы. С. к. е. G является замкнутой нормальной подгруппой в $G$; смежные классы по этой подгруппе совпадают со связными комшонентами группы $G$. Факторгруппа $G / G^{\circ}$ вполне несвязна и хаусдорфова, причем $G^{\circ}$ - наименьшая из таких нормальных подгрупп $H \subset G$, что $G / H$ вполне несвязна. Если $G$ локально связна (напр., $G$ - группа Ли), то $G^{\circ}$ открыта в $G$ и $G / G^{\circ}$ дискретна.



В произвольной алгебраич. группе $G$ С к.е. $G^{\circ}$ также открыта и имеет конечный индекс, причем $G^{\circ}$ является минимальной замкнутой подгруппой конечного индекса в $G$. Связные компоненты алгебраич. группы $G$ совпадают с неприводимыми компонентами. Для любого регулярного гомоморфизма $\varphi: G \rightarrow G^{\prime}$ алгебраич. групп справедливо равенство $\varphi\left(G^{\circ}\right)=\varphi(G)^{\circ}$. Если $G$ огрецелена над нек-рым полем $k$, то и $G^{\circ}$ определена над $k$.



Если $G$ - алгебраич. группа над полем $\mathbb{C}$, то ее $\mathrm{C}$. к.е. $G^{\circ}$ совпадает со С. к. е. группы $G$, рассматриваемой как комплексная группа Ли. Если $G$ определена над $\mathbb{R}$, то группа $G^{\circ}(\mathbb{R})$ вещественных точек в $G^{\circ}$ не обязательно связна в топологии группы Ли $G(\mathbb{R})$, но число ее связных компонент конечно и имеет вид $2^{l}$, где $l \geqslant 0$. Напр., группа $\mathrm{GL}_{n}(\mathbb{C})$ распадается на две компоненты, хотя $\mathrm{GL}_{n}(\mathbb{R})$ связна. Псевдоортогональная унимодуля рная групша $\operatorname{SO}(p, q)$, к-рая может рассматриваться как группа вещественных точек связной комплексной алгебраич. группы $\mathrm{SO}_{p+q}(\mathbb{C})$, связна при $p=0$ или $q=0$ и распадается на две компоненты при $p, q>0$.



Лит.: [1] $\mathbf{6}$ о $\mathrm{p}$ е л ь А.. Линейные алгебраические группы, пер. с англ., М., 1972; [2] П о н т р я г и н Л. С., Непре-

рывные группы, 3 изд., М., 1973; [3] Х е л г а с о н С.. Дифференциальная геометрия и симметрические пространства, пер. с англ., М., 1964; [4] ШI а ф а $\mathrm{p}$ е в и ч И. Р., Основы алгебраической геометрии, М., $1972 . \quad$ А. Л. Онищик.



СВЯЗНАЯ СУММА с е м е й с т в а м н о ж е с т вобъединение этих множеств в единое связное множество. Само понятие С. с. возникло из необходимости отличить такого рода объединение от понятия несвязной или открыто-замкнутой суммы, т. е. такого объединения множеств, когда они не пересекаются и связными подмножествами в этом объединении могут быть только под$\begin{array}{ll}\text { множества-слагаемые. } & \text { В. И. Малыхин. }\end{array}$



СВЯЗНОЕ МНОЖЕСТВО - подмножество объемлющего множества, в к-ром определено понятие связности и в смысле к-рого само подмножество связно. Напр. C. м. пространства действительных чисел являются выпуклые множества и только они; С. м. графа является такое множество, в к-ром любые две точки соединены путем, целиком лежащим в этом множестве.



С. И. Малыхин. СВЯЗНОЕ ПРОСТРАНСТВО - топологическое пространство, к-рое нельзя представить в виде суммы двух отделенных друг от друга частей или, более строго, непустых непересекающихся открыто-замкнутых подмножеств. Пространство связно тогда и только тогда, когда каждая непрерывная числовая функция принимает на нем все промежуточные значения. Непрерывныї образ С. П., произведение С. П., пространство замкнутых подмножеств в топологии Вьеториса С. п. суть С. П. Каждое связное вполне регулярное пространство имеет мощность не менее континуума, хотя существуют и счетные связные хаусдорфовы пространства.



СВЯЗНОСТ В. И. Малыхин. СВЯЗНОСТИ НА МНОГООБРАЗИИ - дифференциально-геометрические структуры на гладком многообразии $M$, являющиеся связностяли в приклеенных к $M$ гладких расслоенных пространствах $E$ с однородными типовыми слоями $G / H$ размерности $\operatorname{dim} M$. В зависимости от выбора однородного пространства $G / H$ получаются, напр., аббинные связности, проективные связнасти, конборлные связности и др. на многообразии $M$. Общее понятие С. на м. ввел Э. Картан [1]; он назвал многообразие $M$ с заданной на нем связностью «неголономным пространством с фуундаментальной группой».



Современное определение С. на м. опирается на понятие гладкого расслоенного пространства, приклеенного к многообразию $M$. Пусть $F=G / H$ является однородныл пространством размерности $\operatorname{dim} M$ (напр., аффқинным пространством, проективным пространством и т. п.). Пусть $p: E \rightarrow M$ является гладким локально тривиальным расслоением с типовым слоем $F$ и пусть в әтом расслоении существует и фиксировано гладкое сечение $s$, т. е. такое гладкое отображение $s: M \rightarrow E$, что $p(s(x))=x$ для любого $x \in M$. Последнее условие гарантирует, что $s$ является диффеоморфизмом $M$ на $s(M)$, и поэтому $M$ и $s(M)$ можно при желании отождествлять. Другими словами, к каждой точке $x \in M$ присоединен экземпляр $F_{x}$ однородного пространства $F$ размерности $\operatorname{dim} M$ - слой расслоения $p: E \rightarrow M$ над $x-c$ фиксированной в ней точкой $s(x)$, отождествляемой с $x$.





\end{document}